\begin{corollary}[E.25 --- Proved: the remaining rational class lies in an exact five-dimensional subspace]
\label{E-five-parameters}
Let $A^+$ be the integral elliptic-involution submodule of
E.24 and retain $D$ and $\mathcal Q$ from E.23. Let
$\mathcal Q^+$ and $V^+$ denote the same involution's fixed
submodules. The actual rational class $\eta$ belongs to the
inverse image
\[
 \mathcal H=\operatorname{preimage}\left(
 H^1(\mathcal G,D(-1))\subset
 H^1(\mathcal G,\mathcal Q^+_{\mathbb Q_3}(-1))\right)
\]
under the specified map from
$H^1(\mathcal G,A^+_{\mathbb Q_3}(-1))$. There is a
short exact sequence of rational vector spaces
\[
 0\longrightarrow H^1(\mathcal G,V^+_{\mathbb Q_3})
   \longrightarrow\mathcal H
   \longrightarrow H^1(\mathcal G,D(-1))
   \longrightarrow0,
 \qquad \dim\mathcal H=3+2=5.
 \tag{E25.1}
\]
Thus the rational action class remaining after E.19 is
confined to this actual five-dimensional subspace. The
integral torsion restriction remains the separate one
computed in E.24. The five coordinates are not evaluated
by this dimension calculation.
\end{corollary}
\begin{proof}
The involution fixes $w$, so the equivariant isomorphism
$B_w(-1)\simeq V$ of E.23 identifies
$B_w^+(-1)$ with $V^+$. Its projector $(1+\iota)/2$ is
integral and exact, giving
\[
 0\longrightarrow V^+_{\mathbb Q_3}
   \longrightarrow A^+_{\mathbb Q_3}(-1)
   \longrightarrow\mathcal Q^+_{\mathbb Q_3}(-1)
   \longrightarrow0.
\]
The right-hand representation has zero invariants, since
its good-Frobenius eigenvalues have positive weights.
The left-hand $H^2$ is zero by the unramified arithmetic
source calculation of E.16. Its cohomology sequence is
therefore short exact in degree one.

On $V$, the elliptic involution has compact trace $-2$
and boundary trace zero; hence $\dim V^+=4$.
Complex conjugation and its product with that involution
both have trace $-2$ on $V$: on the cusp set each fixes
three points, with the boundary Tate sign, while its
compact trace is zero. Thus
\[
 \operatorname{tr}(c|V^+)=\frac{-2-2}{2}=-2.
\]
The same Euler-characteristic calculation, with zero
$H^0$ and $H^2$, gives $\dim H^1(\mathcal G,V^+)=3$.
Because $\iota$ is an element of $\mathcal A$, the
subspace $D$ lies in $\mathcal Q^+_{\mathbb Q_3}$.
The $\mathcal A$-averaging projector commutes with Galois
and with $\iota$; hence its inclusion is split and so
is injective on $H^1$. Pulling back this two-dimensional
subspace through the short exact cohomology sequence
proves (E25.1). E.24 puts the actual $\eta$ in the source,
and E.23 puts its image in the specified subspace, so it
indeed belongs to $\mathcal H$. The independent finite
character program records the two traces used here.
All operations are the displayed projector, bracket,
quotient, and inverse-image maps; no replacement of the
arithmetic class by an unspecified graded action occurs.
The anchor is again \S51.16, folios 733v--736v, scans
172--178, and LM7-E.23(i).
\end{proof}
